\subsection{Guarantees}\label{sec:theory}

Let $Q$ be the law of a variant obtained by drawing a residue from the gene's population of assayed residues and then
one of its variants uniformly at random (binary-labelled, non-splice variants outside the ClinVar labelled set). Every
rate and likelihood ratio below is taken under $Q$, so it weights residues equally.

\begin{assumption}[Clustered sampling]\label{as:cluster}
Residues in $\mathcal A$, in $\mathcal B$ and of a test variant are independent draws from one residue population;
within a residue, the kept or tested variant is chosen uniformly at random; $q$ does not depend on these data.
\end{assumption}

\begin{assumption}[Class-conditional invariance]\label{as:shift}
Target variants have, for $k\in\{0,1\}$, the same law of $q$ given $Y=k$ as a test variant under
Assumption~\ref{as:cluster}; their LoF share $\pi=\Pr(Y=1)$ may differ.
\end{assumption}

\begin{proposition}[Per-class coverage]\label{prop:mondrian}
Under Assumption~\ref{as:cluster}, for a fixed $\alpha_k$, with the in-expectation rank and conditionally on $n_k$,
$1-\alpha_k\le\Pr(k\in C\mid Y=k,n_k)\le 1-\alpha_k+1/(n_k+1)$, the upper bound holding when there are no ties.
\end{proposition}

\begin{proposition}[High-probability per-class coverage]\label{prop:pac}
Under Assumption~\ref{as:cluster}, for a fixed $\alpha_k$ and with the PAC rank, $\Pr(k\in C\mid Y=k,\hat t_k)\ge1-\alpha_k$
with probability at least $1-\delta_{\mathrm P}$ over the calibration data for each $k$, and for both classes with
probability at least $1-2\delta_{\mathrm P}$.
\end{proposition}

\begin{proposition}[Coverage after selection]\label{prop:select}
Under Assumption~\ref{as:cluster}, for any $\hat\alpha\in G$ chosen from data independent of the test variant
(for example by \cref{alg:certivar}), the in-expectation rank gives $\Pr(k\in C\mid Y=k)\ge1-\max G$; the PAC rank at
level $\delta_{\mathrm P}/|G|$ gives $\Pr(k\in C\mid Y=k,\hat t_k)\ge1-\hat\alpha$ for each $k$ with probability at least
$1-\delta_{\mathrm P}$.
\end{proposition}

\begin{theorem}[Certified likelihood ratios]\label{thm:lr}
Under Assumption~\ref{as:cluster}, conditionally on $\mathcal A$ and with probability at least $1-\delta$ over
$\mathcal B$, $\mathrm{LR}^+(\hat t(\alpha))\ge L^+(\alpha)$ for every $\alpha\in G$ simultaneously; the same holds for
$\mathrm{LR}^-\le U^-$ with a second $\delta$. Because the event does not depend on the level, each call type returned
by \cref{alg:certivar} reaches its certified level at every evidence level reported, with probability at least
$1-\delta$ (both call types with probability at least $1-2\delta$), whichever $\alpha$ is selected.
\end{theorem}

\begin{theorem}[Fixed-sequence certification]\label{thm:fs}
Under Assumption~\ref{as:cluster}, fix a level $\tau$, a confidence $\delta'$ and an order of $G$ chosen in advance.
Reject $H_\alpha:\mathrm{LR}^+(\hat t(\alpha))\le\tau$ when $\underline{\lambda}_1/\overline{\lambda}_0>\tau$ with both
endpoints at level $\delta'/2$, stop at the first $\alpha$ not rejected, and select the rejected $\alpha$ with the
largest yield; the rejected set is a prefix of the order. (a) Conditionally on $\mathcal A$, the selected LoF calls have
$\mathrm{LR}^+>\tau$ with probability at least $1-\delta'$. (b) If each of $L$ levels is tested in this way with
$\delta'=\delta/L$, every level claimed for LoF calls holds simultaneously with probability at least $1-\delta$. The same
holds for functional calls with $\overline{\phi}_1/\underline{\phi}_0<\tau'$, and for both call types together with
probability at least $1-2\delta$.
\end{theorem}

\begin{corollary}[No prior needed]\label{cor:shift}
Under Assumptions~\ref{as:cluster} and~\ref{as:shift}, for target variants drawn residue first as under $Q$,
Propositions~\ref{prop:mondrian}--\ref{prop:select} and Theorems~\ref{thm:lr}--\ref{thm:fs} hold whatever their LoF
share $\pi\in(0,1)$: for a fixed
$\alpha$, LoF variants are called functional at rate at most $\alpha_1$, and the certified likelihood ratios with
respect to the assay outcome need no prior. For any $\alpha^+,\alpha^-\in G$, the LoF calls at $\alpha^+$ and the
functional calls at $\alpha^-$ are disjoint.
\end{corollary}

Propositions~\ref{prop:mondrian}--\ref{prop:pac} are the class-conditional and training-conditional conformal
guarantees \cite{Vovk2005,Sadinle2019,Vovk2013,Tong2018} under $Q$; Theorems~\ref{thm:lr}--\ref{thm:fs} apply the union
bound of risk control \cite{Bates2021rcps,Geifman2017} and the fixed-sequence argument of Learn-then-Test
\cite{Angelopoulos2025ltt} to a \emph{ratio} of class-conditional rates of an abstaining rule. Monotonicity of both ranks
in $\alpha$ makes coverage at a selected $\alpha$ at least that at $\max G$ and keeps the call types disjoint. Proofs are
in \cref{app:proofs}.

\paragraph{Scope.}
After selection, per-class coverage is guaranteed only at $1-\max G$; Theorem~\ref{thm:fs} certifies levels, not
bounds. Assumption~\ref{as:cluster} treats residues as independent draws (our random splits draw without replacement;
the oracle design does not), and a contiguous region need not be exchangeable with the rest; Assumption~\ref{as:shift}
holds only approximately for VUS (\cref{sec:shift}) and not for variants whose assay result would be intermediate.
